\section*{Hoofdstuk \ref{ch:b2:ligand-field} --- Ligandveldtheorie en kleur}

\begin{solution}{exo:b2:ligand-field:1}
$d^3$: $t_{2g}^3$, CFSE $-1.2\Delta_o$. $d^8$: $t_{2g}^6e_g^2$, CFSE $-2.4 + 1.2 = -1.2\Delta_o$.
\end{solution}

\begin{solution}{exo:b2:ligand-field:2}
\ce{[Fe(H2O)6]^2+}, $d^6$ \omterm{def:b2:ligand-field:spin}{hoogspin} $t_{2g}^4e_g^2$: vier. \ce{[Fe(CN)6]^4-}, \omterm{def:b2:ligand-field:spin}{laagspin}
$t_{2g}^6$: geen (\omterm{def:b2:diatomic-mos:magnetism}{diamagnetisch}).
\end{solution}

\begin{solution}{exo:b2:ligand-field:3}
Licht van 600 nm is oranje; het complex ziet er blauw uit (groenachtig blauw).
\end{solution}

\begin{solution}{exo:b2:ligand-field:4}
$\ce{Cl-} < \ce{H2O} < \ce{NH3} < \ce{CN-}$.
\end{solution}

\begin{solution}{exo:b2:ligand-field:5}
Water: $\Delta_o < P$, \omterm{def:b2:ligand-field:spin}{hoogspin}, $t_{2g}^3e_g^2$, vijf ongepaarde elektronen. Cyanide: $\Delta_o > P$,
\omterm{def:b2:ligand-field:spin}{laagspin}, $t_{2g}^5$, één ongepaard elektron.
\end{solution}

\begin{solution}{exo:b2:ligand-field:6}
In het tetraëdrische \ce{[CoCl4]^2-} is de splitsing ongeveer $4/9$ van een octaëdrische, en
chloride is een zwakker ligand dan water: de d-d-band verschuift naar langere golflengten (van oranje
naar rood), en het complex ziet er blauw uit. Tetraëdrische complexen absorberen ook sterker,
vandaar de diepe kleur.
\end{solution}

\begin{solution}{exo:b2:ligand-field:7}
\ce{Ni^2+} is $d^8$. Vlak vierkant: vier lagere orbitalen bevatten de acht elektronen in paren,
$\mathrm d_{x^2-y^2}$ blijft leeg: \omterm{def:b2:diatomic-mos:magnetism}{diamagnetisch}. Tetraëdrisch: $e^4t_2^4$, twee ongepaarde elektronen
in de verzameling $t_2$.
\end{solution}

\begin{solution}{exo:b2:ligand-field:8}
$\tilde\nu = 1/(500 \times 10^{-7}\,\text{cm}) = \qty{20000}{cm^{-1}}$; $\Delta_o =
0.1196/(500 \times 10^{-9}) = \qty{2.39e5}{J/mol}$, \qty{239}{kJ/mol}.
\end{solution}

\begin{solution}{exo:b2:ligand-field:9}
\ce{Zn^2+} is $d^{10}$ en \ce{Sc^3+} $d^0$: geen \omterm{def:b2:ligand-field:d-d}{d-d-overgang} is mogelijk (geen leeg d-niveau
voor \ce{Zn^2+}, geen d-elektron voor \ce{Sc^3+}), dus geen absorptie in het zichtbare.
\end{solution}

\begin{solution}{exo:b2:ligand-field:10}
Co(III) $d^6$ + zes vrije elektronenparen van ammoniak: 18 elektronen. Twaalf vullen de zes \omterm{def:b2:diatomic-mos:bonding}{bindende orbitalen},
zes vullen het niet-bindende $t_{2g}$ (ammoniak geeft met kobalt(III) een veld dat sterk genoeg is voor
\omterm{def:b2:ligand-field:spin}{laagspin}); $e_g^*$ is leeg. Alle elektronen zijn gepaard: \omterm{def:b2:diatomic-mos:magnetism}{diamagnetisch}.
\end{solution}

\begin{solution}{exo:b2:ligand-field:11}
\ce{CO} is een goede $\sigma$-donor en vooral een $\pi$-acceptor: zijn lege $\pi^*$-orbitalen
verlagen het niveau $t_{2g}$ door \omterm{def:b2:ligand-field:pi-ligands}{terugdonatie}, wat $\Delta_o$ vergroot. \ce{F-} is een $\pi$-donor, die
$t_{2g}$ verhoogt en $\Delta_o$ verkleint. Het puntladingsmodel ziet alleen de lading en krijgt de
volgorde verkeerd.
\end{solution}

\begin{solution}{exo:b2:ligand-field:12}
Zonder veld zouden de hydratatie-enthalpieën een gladde lijn volgen (de ionen krimpen langs
de reeks). Water geeft hoogspin-aquacomplexen waarvan de CFSE nul is voor $d^0$, $d^5$, $d^{10}$
en het grootst voor $d^3$ en $d^8$: de extra stabilisatie voegt twee bulten toe, met maxima rond
\ce{V^2+} en \ce{Ni^2+}, boven de lijn door \ce{Ca^2+}, \ce{Mn^2+}, \ce{Zn^2+}.
\end{solution}

\begin{solution}{pb:b2:ligand-field:1}
\textbf{1.} $d^3$.
\textbf{2.} $t_{2g}^3$, één elektron in elk van de drie orbitalen, met parallelle spins.
\textbf{3.} Nee: drie elektronen passen in $t_{2g}$ zonder paring.
\textbf{4.} $-1.2\Delta_o$.
\textbf{5.} Drie.
\textbf{6.} De CFSE is groot en de $e_g^*$-orbitalen, die naar de liganden wijzen, zijn
leeg: een ligand weghalen of toevoegen kost een groot deel van die stabilisatie.
\textbf{7.} Zes oxide-ionen op de hoekpunten van een licht vervormde octaëder.
\textbf{8.} $1/(556 \times 10^{-7}) = \qty{17990}{cm^{-1}}$, ongeveer \qty{1.80e4}{cm^{-1}}.
\textbf{9.} $0.1196/(556 \times 10^{-9}) = \qty{2.15e5}{J/mol}$: \qty{215}{kJ/mol}.
\textbf{10.} Geelgroen (556 nm) en violet (405 nm).
\textbf{11.} Rood, met een beetje blauw: het diepe rood van robijn.
\textbf{12.} Een $d^3$-ion heeft meer dan één aangeslagen configuratie met een elektron in $e_g$;
de elektronenafstoting geeft ze verschillende energieën, vandaar verscheidene banden (behandeld in het
deel van jaar~3).
\textbf{13.} $0.1196/(620 \times 10^{-9}) = \qty{1.93e5}{J/mol}$, \qty{193}{kJ/mol}
(\qty{16130}{cm^{-1}}).
\textbf{14.} Robijn: zijn splitsing is groter.
\textbf{15.} Naar oranjerood: nu wordt rood geabsorbeerd, en groen (met wat blauw) wordt doorgelaten.
\textbf{16.} In beril zijn de chroomplaatsen iets groter en de oxide-ionen verder weg,
en de omgeving verschilt: een zwakker veld.
\textbf{17.} Nee: een $d^3$-ion houdt in elk octaëdrisch veld drie ongepaarde elektronen.
\textbf{18.} $d^9$.
\textbf{19.} Het aquacomplex absorbeert in het rood en oranje door een \omterm{def:b2:ligand-field:d-d}{d-d-overgang}: het
doorgelaten licht is blauw.
\textbf{20.} Het wordt wit: zonder waterliganden rond het koper (sulfaat neemt hun
plaats in) is het veld zwakker en verschuift de band naar het infrarood.
\textbf{21.} Water bindt opnieuw aan koper en het blauwe aquacomplex ontstaat.
\textbf{22.} Bij kortere golflengte: ammoniak, hoger dan water in de reeks, geeft een grotere
splitsing (het diepe blauwviolet van het amminecomplex).
\textbf{23.} Nee: $t_{2g}^6e_g^3$ is de enige configuratie.
\textbf{24.} Ja: één ongepaard elektron.
\textbf{25.} $\boldsymbol{\Delta_o \approx \qty{1.8e4}{cm^{-1}}}$, ongeveer
$\boldsymbol{\qty{215}{kJ/mol}}$.
\end{solution}
